\begingroup
\setlength{\abovedisplayskip}{6pt}
\setlength{\belowdisplayskip}{6pt}
\setlength{\abovedisplayshortskip}{3pt}
\setlength{\belowdisplayshortskip}{4pt}

% Printed 102; principal PDF 109; Google PDF 113.
\par\clearpage
\noindent\makebox[12mm][l]{\fontsize{13}{16}\selectfont 102}\makebox[\dimexpr\linewidth-24mm\relax][c]{\footnotesize\textbf{BRIEFWECHSEL.\quad ÄNDERUNG DER POLHÖHE MIT DER HÖHE.}}\makebox[12mm][r]{}\par\medskip

\[
\theta=\frac{g'-g}{g}=\tfrac54\frac hg
\]
\noindent gilt (die \textsc{Newton}sche Abplattung).

Bei dieser neuen umfassenden Form wird, wenn ich mit \textsc{Sabine}
\[
\frac{g'-g}{g}=\frac1{192{,}7}
\]
\noindent setze, die Polhöhe in \(a\)
\[
=\varphi+1070''\,.\,\frac{s}{a_0}\sin 2\varphi.
\]

Für die höchsten Berge in Schlesien also nur etwa \(\tfrac14\) Sekunde. Ich muss
Ihnen indessen offen gestehen, dass ich die ganze Untersuchung nur als
eine theoretische Kuriosität betrachten kann, der jede praktische Bedeutung
fehlt. Sie hätte eine solche nur dann, wenn auf der glatten Erdoberfläche
\(DAE\) eine dünne, hohe Säule \(Aa\) errichtet wäre, an deren Spitze
wie an ihrem Fuß man die Polhöhe beobachten könnte. In Wirklichkeit, wenn
\(a\) etwa auf einem hohen Berg liegt, ist zunächst einmal der Punkt \(A\) gar
nicht zugänglich. Und selbst wenn man ihn erreichen und den Unterschied in der
Richtung der Lotlinie in \(A\) und \(a\) durch Messungen genau bestimmen
könnte, wäre man doch nicht berechtigt, die obige Formel als Darstellung
dieses Unterschieds anzusehen, da die Anziehungskräfte der oberhalb des
Niveaus von \(A\) liegenden Bestandteile des Erdkörpers viel größere und
rechnerisch nicht erfassbare Unterschiede in den Endergebnissen
für die Schwere in \(A\) und \(a\) hervorrufen werden.

Ich habe mich zu diesem Gegenstand im Allgemeinen in meiner Schrift
von 1828 über den Breitenunterschied von Göttingen und Altona, S.~73 [vgl.\
diesen Band, S.~49], bereits so geäußert, dass ich jetzt nichts Besseres
darüber zu sagen weiß.

\begin{center}\rule{.18\textwidth}{.4pt}\end{center}
{\small
\begin{center}\textbf{BEMERKUNGEN.}\end{center}
Der vorstehende Brief, von dem eine Abschrift im Gauss-Archiv vorliegt, ist abgedruckt in den
»Protocollen der Verhandlungen der permanenten Commission der europäischen Gradmessung vom 23. bis
29. September 1869 in Florenz«. (Als Manuskript gedruckt.) S.~30---32.
\par\hfill\textsc{Krüger, Börsch.}\hspace{.06\textwidth}\par
}
\begin{center}\rule{.24\textwidth}{.4pt}\end{center}

